\chapter{Red Blood Cell Antibody Screen}

\section*{What Is This Test?}
The red blood cell antibody screen, also known as the indirect antiglobulin test (IAT) or indirect Coombs test, detects the presence of unexpected, clinically significant antibodies in a patient's serum that react against non-self (allogeneic) red blood cell antigens. These alloantibodies are typically IgG immunoglobulins that develop following exposure to foreign red blood cell antigens through pregnancy (fetomaternal hemorrhage), blood transfusion, or solid organ/hematopoietic stem cell transplantation. The antibody screen uses commercially available group O reagent red blood cells (screening cells) that carry a comprehensive panel of clinically significant antigens, including those of the Rh, Kell, Duffy, Kidd, MNS, and Lewis blood group systems. The patient's serum is incubated with these screening cells, and if antibodies are present, they bind to the corresponding antigens. After incubation, the cells are washed to remove unbound immunoglobulins, and anti-human globulin (AHG, Coombs reagent) is added. The AHG bridges antibody-coated red blood cells, causing visible agglutination. A positive antibody screen requires further identification using a panel of 11--16 reagent red blood cells with known antigen profiles to determine the specificity of the alloantibody. The antibody screen is a critical pre-transfusion test, essential for preventing hemolytic transfusion reactions, and is also used in prenatal care to detect maternal antibodies that can cause hemolytic disease of the fetus and newborn (HDFN).

\section*{Alternative Names}
RBC Antibody Screen; Indirect Antiglobulin Test (IAT); Indirect Coombs Test; Antibody Detection Test; Type and Screen (T\&S); Prenatal Antibody Screen; Alloantibody Screen; Unexpected Antibody Test

\section*{Principle}
The red blood cell antibody screen is performed using the indirect antiglobulin test. The patient's serum (or plasma) is incubated with group O reagent screening red blood cells from 2 or 3 individual donors at 37\textdegree{}C for 30--60 minutes. The screening cells are selected to express a comprehensive range of clinically significant blood group antigens, including D, C, c, E, e (Rh system), K, k (Kell), Fy\textsuperscript{a}, Fy\textsuperscript{b} (Duffy), Jk\textsuperscript{a}, Jk\textsuperscript{b} (Kidd), M, N, S, s (MNS), Le\textsuperscript{a}, Le\textsuperscript{b} (Lewis), and P\textsubscript{1}. If antibodies in the patient's serum react with corresponding antigens on the screening cells, antibody binding occurs. After incubation, the red cells are washed thoroughly with saline to remove unbound serum proteins (which would neutralize the AHG reagent). Anti-human globulin reagent (Coombs serum), which contains antibodies against human IgG and complement (C3d), is added. If the patient's antibodies are attached to the red cells, the AHG bridges the coated cells, producing visible agglutination. Enhancement media -- low-ionic-strength solution (LISS), polyethylene glycol (PEG), or albumin -- are used to reduce the zeta potential between red cells, accelerating antibody binding and increasing sensitivity. Current practice commonly employs gel column technology (Ortho ID-MTS, Grifols DG Gel) or solid-phase red cell adherence assays (Immucor Capture), which provide standardized, objective, and highly sensitive detection. If the screen is positive, antibody identification is performed using a panel of 11--16 group O reagent red blood cells with known antigen phenotypes, tested against the patient's serum to identify the antibody specificity. Autocontrol (patient's serum tested against patient's own red cells) helps distinguish alloantibodies from autoantibodies.

\section*{History}
The Coombs test was developed by Robin Coombs, Arthur Mourant, and Robert Race at the University of Cambridge, first published in 1945. Their groundbreaking paper demonstrated that rabbit anti-human globulin serum could detect non-agglutinating (incomplete) Rh antibodies bound to red blood cells. Before the Coombs test, only IgM antibodies that directly agglutinated red cells in saline (complete antibodies) could be detected, leaving many clinically important IgG antibodies undetectable. The direct antiglobulin test (DAT, direct Coombs) detects antibodies already bound to red cells in vivo and is used to diagnose autoimmune hemolytic anemia, HDFN, and drug-induced hemolysis. The indirect antiglobulin test (IAT, indirect Coombs), which is the antibody screen, detects free antibodies in the serum. The IAT became a standard pre-transfusion test in the 1950s--1960s, dramatically reducing the incidence of hemolytic transfusion reactions. Gel column technology (Lapierre, 1990) and solid-phase red cell adherence assays (1990s) introduced automation and standardization to the historically manual tube method. Today, all pre-transfusion testing in developed countries includes both ABO/Rh typing and an antibody screen (type and screen).

\section*{Why This Test Is Ordered}
\begin{itemize}
  \item Pre-transfusion testing --- all patients receiving blood products require an antibody screen as part of the ``type and screen'' to detect alloantibodies that could cause hemolytic transfusion reactions. If the screen is positive, the antibody must be identified and antigen-negative donor units must be selected and crossmatched.
  \item Prenatal testing --- all pregnant women should have an ABO/Rh type and antibody screen at the first prenatal visit. A positive screen identifies maternal alloantibodies (most importantly anti-D, but also anti-Kell, anti-c, anti-E) that can cross the placenta and cause hemolytic disease of the fetus and newborn (HDFN). If positive, antibody titers are monitored serially.
  \item RhD-negative pregnant women --- antibody screen before administration of Rh immune globulin (RhoGAM) at 28 weeks and postpartum. A negative screen is required before administration.
  \item Evaluation of transfusion reactions --- in a hemolytic transfusion reaction, the antibody screen may be positive and crossmatch may be incompatible
  \item Preoperative evaluation --- type and screen is ordered before surgeries with potential for significant blood loss to reserve compatible blood
  \item Post-transfusion --- repeat type and screen after recent transfusion to detect new alloantibody formation
  \item Evaluation of unexplained anemia --- the DAT differentiates immune from non-immune hemolysis, and the IAT detects circulating alloantibodies or autoantibodies
\end{itemize}

\section*{Is This a Primary or Secondary Test}
The RBC antibody screen is a primary pre-transfusion test required for all patients receiving blood products and a standard prenatal screening test for all pregnant women.

\section*{How This Test Is Performed}
A venous blood sample is collected in a plain red-top tube (serum, no anticoagulant) or a pink/lavender-top (EDTA) tube. For serum, approximately 5--7 mL of blood is drawn and allowed to clot, then centrifuged to separate serum. For plasma, EDTA-anticoagulated whole blood is centrifuged to obtain plasma. The serum or plasma is incubated with 2--3 group O reagent screening red blood cells of known antigen profile at 37\textdegree{}C with enhancement medium (LISS, PEG, or albumin) for 30--60 minutes. After incubation, the cells are washed 3--4 times with saline, and AHG reagent is added. Following centrifugation, agglutination is assessed visually (tube method), by gel column filtration (gel method --- agglutinated cells are trapped at the top of the gel column), or by automated solid-phase reading (solid-phase method --- bound cells remain after washing, detected by indicator cells). A positive result (visible agglutination) indicates that antibodies in the patient's serum recognize antigens on the screening cells. The antibody identification panel is then performed using 11--16 additional panel cells with known phenotypes to determine the specific antigen(s) targeted by the antibody. The entire testing process takes 45--90 minutes for the screen, plus additional time for antibody identification if positive.

\section*{What Preparation Is Needed}
No special preparation is required. Fasting is not necessary. The patient should inform the healthcare provider of all medications, with particular attention to recent IVIG administration (passive transfer of antibodies), RhoGAM administration within the past 3 months, recent blood transfusion (may introduce donor antibodies or alter antibody detection), and certain drugs that can cause positive antibody screens (methyldopa, penicillin, cephalosporins, procainamide, quinidine, and many others that cause drug-induced antibodies). Recent or current pregnancy should be disclosed. There is no need to discontinue any medications before the test.

\section*{Contraindications and Precautions}
No absolute contraindications to testing. Important technical considerations include: the antibody screen detects only IgG antibodies (incomplete antibodies) by the IAT method, which require the AHG reagent for visualization --- IgM antibodies (which agglutinate directly in saline at room temperature, such as anti-A, anti-B, some anti-M, anti-N, anti-Le\textsuperscript{a}, anti-Le\textsuperscript{b}, and anti-P\textsubscript{1}) may be missed if the incubation is performed only at 37\textdegree{}C and the immediate spin phase is omitted; some IgM antibodies are cold-reactive and may be clinically insignificant at body temperature, but anti-Le\textsuperscript{a} and anti-P\textsubscript{1} can rarely cause hemolytic transfusion reactions; the antibody screen is performed at 37\textdegree{}C to detect clinically significant antibodies; a positive DAT (red cells already coated with IgG in vivo) does not affect the IAT (which tests serum); recent transfusion (<3 months) can introduce donor antibodies into the patient's serum, potentially causing a positive screen that may not represent the patient's permanent antibody status; recent IVIG administration, RhoGAM (anti-D immune globulin), or monoclonal antibody therapies (daratumumab for multiple myeloma --- anti-CD38 antibody that interferes with all pre-transfusion testing) can cause positive antibody screens that complicate transfusion compatibility testing; for patients on daratumumab, dithiothreitol (DTT) treatment of reagent cells is used to denature CD38 and eliminate the interference; severely lipemic, hemolyzed, or icteric samples may interfere with visual or automated agglutination detection; the sample for pre-transfusion testing must be collected within 3 days of planned transfusion in patients who have been pregnant or transfused within the past 3 months, as new antibody formation can occur rapidly.

\section*{Risks}
Risks are those of routine venipuncture: mild pain, bruising, hematoma, lightheadedness, and rarely, infection or nerve injury. No additional risks are associated with the antibody detection test itself.

\section*{What Happens After the Test}
Results are typically available within 1--2 hours for the antibody screen. A negative antibody screen (no unexpected antibodies detected) means the patient can receive ABO/Rh-compatible blood after an immediate-spin or electronic crossmatch in most institutions. If the screen is positive at any time, the antibody must be identified to its specificity, and any future transfusions require crossmatch-compatible, antigen-negative donor units. Positive prenatal antibody screens require: identification of the antibody specificity, assessment of its potential to cause HDFN (anti-D, anti-c, anti-Kell, anti-E, and anti-Kidd are notable for causing moderate-to-severe HDFN; anti-P\textsubscript{1}, anti-Le\textsuperscript{a}, anti-Le\textsuperscript{b} are generally IgM and do not cross the placenta), serial antibody titers (every 2--4 weeks if the antibody is clinically significant for HDFN), and fetal monitoring with middle cerebral artery Doppler ultrasound if titers reach a critical threshold (indicating significant fetal anemia). A positive antibody screen in a pre-transfusion patient delays transfusion, as antigen-negative units must be found (which may take hours to days for rare antibody combinations or patients with multiple alloantibodies).

\section*{Understanding Results}

\subsection*{Negative (No Unexpected Antibodies Detected)}
\begin{itemize}
  \item \textbf{Normal finding} in individuals who have never been exposed to foreign RBC antigens through transfusion or pregnancy.
  \item \textbf{Indicates} that the patient can safely receive ABO/Rh-compatible blood with standard compatibility testing, assuming no prior history of alloantibodies.
  \item \textbf{Does not guarantee} complete safety; newly formed antibodies following a recent transfusion or pregnancy may not yet be at detectable titers, and rare antibodies against low-frequency antigens not present on screening cells may be missed.
  \item \textbf{In prenatal patients:} A negative screen throughout pregnancy indicates low risk of HDFN; RhD-negative women with a negative screen receive prophylactic RhIG at 28 weeks and after delivery.
\end{itemize}

\subsection*{Positive (Unexpected Antibodies Detected)}
\begin{itemize}
  \item \textbf{Anti-D (Rh system):} The most immunogenic and historically most clinically significant blood group antibody. Can cause severe HDFN (hydrops fetalis) and severe hemolytic transfusion reactions. Incidence has declined dramatically since RhIG prophylaxis was introduced in 1968, but anti-D still occurs in RhD-negative women who did not receive RhIG, in RhD-negative transfusion recipients, and in individuals with partial D phenotype who make anti-D against epitopes they lack.
  \item \textbf{Anti-Kell (anti-K, anti-k, anti-Kp\textsuperscript{a}, anti-Kp\textsuperscript{b}):} Anti-K is the most common Kell system antibody. Causes severe HDFN (suppresses fetal erythropoiesis in addition to immune destruction) and hemolytic transfusion reactions. Kell antibodies are typically IgG and clinically significant.
  \item \textbf{Anti-c and anti-E (Rh system):} Common Rh antibodies in RhD-positive individuals (especially anti-E). May cause HDFN and delayed hemolytic transfusion reactions.
  \item \textbf{Anti-Fy\textsuperscript{a} and anti-Fy\textsuperscript{b} (Duffy system):} Cause delayed hemolytic transfusion reactions and occasionally HDFN. Anti-Fy\textsuperscript{a} is particularly common in multiply transfused patients. Duffy-null phenotype (Fy(a-b-)) is protective against Plasmodium vivax malaria.
  \item \textbf{Anti-Jk\textsuperscript{a} and anti-Jk\textsuperscript{b} (Kidd system):} Notorious for causing severe delayed hemolytic transfusion reactions because antibody titers rapidly fall below detectable levels but rise rapidly upon re-exposure to the antigen, resulting in brisk hemolysis of transfused cells. The classic cause of delayed hemolytic transfusion reactions presenting 5--10 days post-transfusion with falling hemoglobin, jaundice, and hemoglobinuria.
  \item \textbf{Anti-M, anti-N, anti-S, anti-s (MNS system):} Anti-M and anti-N are usually IgM cold-reactive antibodies that are clinically insignificant, but can occasionally be IgG and cause HDFN or transfusion reactions. Anti-S and anti-s are usually IgG, clinically significant, and can cause HDFN and transfusion reactions.
  \item \textbf{Anti-Le\textsuperscript{a} and anti-Le\textsuperscript{b} (Lewis system):} Usually IgM, cold-reactive, and clinically insignificant (do not cause HDFN or hemolytic transfusion reactions because Lewis antigens are not intrinsic to the red cell membrane but are adsorbed from plasma). Rare cases of Lewis antibodies causing hemolytic transfusion reactions have been reported. Lewis antibodies are most commonly encountered in pregnant women (increased incidence, decreased Lewis antigen expression during pregnancy) and are rarely clinically significant.
  \item \textbf{Anti-P\textsubscript{1} (P blood group system):} Usually a cold-reactive IgM antibody, naturally occurring, and clinically insignificant. Rarely IgG; some cases associated with early pregnancy loss.
  \item \textbf{Autoantibodies (warm autoantibodies):} React broadly with all panel cells, including the autocontrol (patient's serum agglutinates the patient's own cells). Associated with warm autoimmune hemolytic anemia (WAIHA), systemic lupus erythematosus, lymphoproliferative disorders (CLL, non-Hodgkin lymphoma), and certain drugs. Identification of underlying alloantibodies is difficult and may require specialized techniques (autoadsorption or alloadsorption).
  \item \textbf{Cold autoantibodies:} IgM autoantibodies reacting at room temperature that may interfere with blood typing and antibody screen, causing false-positive reactions. Associated with cold agglutinin disease, Mycoplasma pneumoniae infection, and Epstein-Barr virus infection. Pre-warming techniques resolve interference.
  \item \textbf{Multiple alloantibodies:} Patients who have received multiple transfusions (sickle cell disease, thalassemia, myelodysplasia) may develop multiple alloantibodies, complicating transfusion support. Extended red cell antigen phenotyping and provision of antigen-matched blood reduce the risk of new alloantibody formation.
  \item \textbf{Antibodies against low-frequency antigens:} Antibodies against antigens not present on screening cells (C\textsuperscript{w}, V, VS, Kp\textsuperscript{a}, Js\textsuperscript{a}, Lu\textsuperscript{a}, Wr\textsuperscript{a}, and many others). The screen is negative, but the antibody may cause crossmatch incompatibility and, rarely, hemolytic transfusion reactions.
\end{itemize}

\section*{Normal Results}
\begin{itemize}
  \item \textbf{Antibody screen:} Negative --- no unexpected, clinically significant red blood cell antibodies detected in the patient's serum.
  \item \textbf{ANTI-D (Rh immune globulin):} Not detected (anti-D from RhoGAM administration is detected by the antibody screen within days to weeks after injection and should not be misinterpreted as active alloimmunization).
  \item \textbf{In prenatal screening:} ABO and RhD type plus antibody screen negative. If the screen is positive at any time during pregnancy, the antibody must be identified, and clinically significant antibodies must have titers determined.
\end{itemize}

\section*{Follow-up Tests}
\begin{itemize}
  \item \textbf{Antibody identification panel:} Performed whenever the antibody screen is positive, using 11--16 panel cells to determine the specific antigen(s) targeted by the antibody
  \item \textbf{Direct antiglobulin test (DAT, direct Coombs test):} To determine if red blood cells are coated with antibody or complement in vivo, for diagnosing autoimmune hemolytic anemia, HDFN, and drug-induced hemolytic anemia
  \item \textbf{Elution:} Removal of antibodies bound to red blood cells to identify their specificity (DAT-positive patients)
  \item \textbf{Red blood cell antigen phenotyping:} Determines which antigens the patient's red cells carry --- essential for identifying compatible blood donors and confirming the absence of an antigen to explain a specific antibody
  \item \textbf{Crossmatch (major crossmatch):} Testing of the patient's serum against specific donor red blood cell units --- must be performed for all patients with a positive antibody screen or history of alloantibodies
  \item \textbf{Serum bilirubin (total and direct), haptoglobin, LDH, reticulocyte count, and peripheral blood smear:} For evaluation of hemolysis when transfusion reaction or autoimmune hemolytic anemia is suspected
  \item \textbf{Antibody titration:} In prenatal patients with clinically significant antibodies (anti-D, anti-c, anti-Kell, anti-E, anti-Kidd), serial titers every 2--4 weeks from 18--20 weeks gestation, with fetal monitoring when critical titers are reached (usually >16--32 depending on antibody specificity)
  \item \textbf{Fetal middle cerebral artery Doppler ultrasound:} Non-invasive assessment of fetal anemia in pregnancies with clinically significant maternal alloantibodies
  \item \textbf{Amniocentesis for delta OD450 spectrophotometry and fetal blood sampling (cordocentesis):} For definitive fetal assessment when non-invasive monitoring suggests significant anemia
  \item \textbf{Cold agglutinin titer:} When cold autoantibodies are detected
  \item \textbf{Donath-Landsteiner antibody test:} For paroxysmal cold hemoglobinuria
  \item \textbf{Drug history review and drug-induced antibody testing:} When drug-induced immune hemolytic anemia is suspected
\end{itemize}

\section*{References}
\begin{enumerate}
  \item Tierney LM, Saint S, Drazen JM. CURRENT Medical Diagnosis \& Treatment. McGraw-Hill; 2023.
  \item Wallach J. Interpretation of Diagnostic Tests. 11th ed. Wolters Kluwer; 2022.
  \item Fischbach FT, Dunning MB. A Manual of Laboratory and Diagnostic Tests. 11th ed. Wolters Kluwer; 2023.
  \item Burtis CA, Ashwood ER, Bruns DE. Tietz Textbook of Clinical Chemistry and Molecular Diagnostics. 7th ed. Elsevier; 2023.
  \item Roback JD, Grossman BJ, Harris T, Hillyer CD, eds. AABB Technical Manual. 20th ed. AABB; 2020.
  \item Fung MK, Eder AF, Spitalnik SL, Westhoff CM, eds. AABB Technical Manual. 19th ed. AABB; 2017.
  \item Klein HG, Anstee DJ. Mollison's Blood Transfusion in Clinical Medicine. 12th ed. Wiley-Blackwell; 2014.
  \item Coombs RRA, Mourant AE, Race RR. A new test for the detection of weak and ``incomplete'' Rh agglutinins. Br J Exp Pathol. 1945;26:255--266.
  \item RCOG. The Management of Women with Red Cell Antibodies during Pregnancy. Green-top Guideline No. 65. Royal College of Obstetricians and Gynaecologists; 2014.
\end{enumerate}

\clearpage
